\section{Implementación mínima realizada}

La prueba funcional se ejecuta en fases numeradas. El esquema crea extensiones,
tablas y restricciones; el seed carga el conjunto sintético; el tercer script
agrega índices y la vista de recuperación; el quinto crea roles, privilegios,
políticas RLS y auditoría. Los archivos de consultas ejecutan los casos C1--C9
y un archivo separado reúne aserciones negativas de integridad y seguridad.

El comando \path{scripts/cargar_base.sh} ofrece un recorrido legible: regenera
los artefactos, recrea la base y muestra los resultados de los casos. El
comando \path{scripts/validar_base.sh} se reserva para la verificación
exhaustiva. Sólo acepta una base cuyo nombre termine en
\texttt{\_validacion}, la recrea dos veces, comprueba checksums y conteos,
ejecuta rechazos controlados, valida permisos y compara snapshots ordenados.

La implementación se detiene allí. No agrega un backend que oculte las reglas
de la base ni una interfaz para simular una aplicación completa. Las pruebas
usan directamente SQL y roles técnicos, que permiten observar si cada contrato
se cumple en la capa de datos.
